\documentclass[12pt, reqno]{article}

\usepackage{amsmath, amsthm, amssymb}
\usepackage{enumitem}
\usepackage{tcolorbox}
\usepackage{hyperref}
\usepackage{tikz}
\usepackage{tikz-cd}
\usepackage{pgfplots}
\pgfplotsset{compat=1.18}
\usetikzlibrary{arrows.meta}
\usepackage{mathrsfs}
\usepackage{fancyhdr}
\usepackage[bottom=0.75in, top=1in, left=0.5in, right=0.5in]{geometry}
\usepackage{array}   % for \newcolumntype macro
\usepackage{lmodern}
\newcolumntype{L}{>{$}l<{$}}


\theoremstyle{plain}
\newtheorem{theorem}{Theorem}[section]
\newtheorem{proposition}{Proposition}[section]
\newtheorem{exercise}{Exercise}[section]
\newtheorem{lemma}{Lemma}[section]
\newtheorem{corollary}{Corollary}[section]

\theoremstyle{definition}
\newtheorem{definition}{Definition}[section]
\newtheorem{example}{Example}[section]

\theoremstyle{remark}
\newtheorem{remark}{Remark}[section]

\renewcommand{\phi}{\varphi}
\renewcommand{\epsilon}{\varepsilon}
\renewcommand{\emptyset}{\varnothing}

\newcommand{\RR}{\mathbb{R}}
\newcommand{\ZZ}{\mathbb{Z}}
\newcommand{\NN}{\mathbb{N}}
\newcommand{\CC}{\mathbb{C}}
\newcommand{\QQ}{\mathbb{Q}}

\DeclareMathOperator{\ima}{\text{im}}

\begin{document}

\topmargin=-40pt
\rhead{Henry Woodburn}
\lhead{Reading Summer 2026}
\renewcommand{\headrulewidth}{1pt}
\renewcommand{\headsep}{20pt}
\thispagestyle{fancy}

{\Huge \bfseries \noindent Roe Notes}

\tableofcontents

\section{Riemannian Geometry}  

\subsection{Connections}

Let $M$ be a smooth manifold, and let $V$ be a vector bundle on $M$. 

\begin{definition}
    A \textit{Connection} on $V$ is a linear map 
    \[
        \nabla: C^\infty(TM)\otimes C^\infty(V) \to C^\infty(V)
    \]
    which assigns to a vector field $X$ and a section $Y$ of $V$ a new section $\nabla_X Y$ 
    of $V$, such that for any smooth function $f$ on $M$,
    \begin{enumerate}
        \item $\nabla_{fX}Y = f\nabla_X Y$
        \item $\nabla_X{fY} = f\nabla_X Y + (X.f)Y$,
    \end{enumerate}
    where $X.f$ denotes the lie derivative of $f$ along $X$. 
\end{definition}

Condition $(1.)$ tells us that $X \mapsto \nabla_X Y$ is a homomorphism of modules over
$C^\infty(M)$. In other words, a connection is tensorial in $C^\infty(TM)$. This means
that we can regard a connection as a map 
\[
    C^\infty(V) \to \Omega^1(V) := C^\infty(T^*M \otimes V),
\]
the space of $V$ valued forms on $M$.

In a local trivialization, an example of a connection is the Lie derivative. Thus many local connections 
exist.

Moreover, the difference of two connections is local in both $X$ and $Y$, making it an $\text{End}(V)$ valued
one-form. Thus, we can use partitions of unity to piece together local connections into a global one.
So connections exist and are plentiful. 

From the above discussion, we deduce that in a local trivialization,
\[
    \nabla_i = \frac{\partial}{\partial x_i} + \Gamma_i,
\]
where $\nabla_i = \nabla_{\frac{\partial}{\partial x_i}}$, and $\Gamma_i$ is 
a smooth section of $\text{End}(V)$. The $\Gamma_i$ are called the Christoffel symbols. 

\begin{remark}
    We can also define connections using \textit{paralell transport}. Let $\gamma$ be a 
    smooth curve in $M$. The differential equation $\nabla_{\dot{\gamma}}Y = 0$ is a first 
    order ODE on sections $Y$ of $V$ along $\gamma$ which has a unique solution for a 
    given initial value. One says that $Y$ is paralell along $\gamma$, or that 
    $Y$ has been obtained from its initial value by paralell transport. 
\end{remark}

\begin{definition}
    The \textit{curvature operator $K$} of a connection $\nabla$ on $V$ is defined as follows: 
    if $X$ and $Y$ are vector fields and $Z$ is a section of $V$, then 
    \[
        K(X,Y)Z = \nabla_X\nabla_Y Z - \nabla_Y \nabla_X Z - \nabla_{[X,Y]}Z.
    \]
\end{definition}

\begin{lemma}
    The curvature $F(X,Y)Z$ at a point $m \in M$ depends only on the values of $X, Y,$ and $Z$ at the point 
    $m$.
\end{lemma}

\textit{Proof:} Choose some local coordinate patch and a trivialization of $V$. Then we can expand
\begin{align}
    \nabla_X \nabla_Y Z - \nabla_Y \nabla_X Z & = 
    \sum_{i,j} \left[X^j \frac{\partial Y^i}{\partial x_j} 
    - Y^j \frac{\partial X^i}{\partial x_j}\right]\frac{\partial Z}{\partial x_i}
    + \left[ X^j Y^i - Y^j X^i\right] \frac{\partial^2 Z}{\partial x_j \partial x_i}\\
    & + \left[ X^j Y^i - Y^j X^i\right] \left(\frac{\partial \Gamma_j(Z)}{\partial x_i} 
    + \frac{\partial\Gamma_i(Z)}{\partial x_j}\right)\\
    & + \left[ X^j \frac{\partial Y^i}{\partial x_j} - Y^j \frac{\partial X^i}{\partial x_j}\right]\left(\Gamma_i(Z)\right)
    + \left[X^j Y^i - Y^j X^i\right]\left(\Gamma_j \Gamma_i(Z)\right).
\end{align}

Since mixed partials commute, the second term in (1) vanishes. By symmetry, the term (2) vanishes as well, leaving only
\begin{align*}
    \nabla_X \nabla_Y Z - \nabla_Y \nabla_X Z & = 
    \sum_{i,j} \left[X^j \frac{\partial Y^i}{\partial x_j} 
    - Y^j \frac{\partial X^i}{\partial x_j}\right]\left(\frac{\partial Z}{\partial x_i} + \Gamma_i(Z)\right)\\
    & + \left[X^j Y^i - Y^j X^i\right]\left(\Gamma_j \Gamma_i(Z)\right)\\
    & = \nabla_{[X,Y]}Z + \left[X^j Y^i - Y^j X^i\right]\Gamma_j\Gamma_i(Z).
\end{align*}

From here the lemma follows immediately. \qed

Thus $K$ is induced by a bundle map $TM \otimes TM \to \text{End}(V)$, and moreover since $K(X,Y)$ is 
antisymmetric in $X$ and $Y$, we may regard $K$ as a 2-form with values in $\text{End}(V)$. In terms of local 
coordinates, we define the curvature 2-form to be the $\text{End}(V)$-valued 2-form 
\[
    K = \sum_{i < j}K(\partial_i, \partial_j)dx^i \wedge dx^j,
\]
where $K(\partial_i, \partial_j)$ takes values in $\text{End}(V)$. 

Now we consider connections on the tangent bundle $TM$. Let $x^i$ be local coorinates with a corresponding
vector fields $\partial_i$. Then we can express the endomorphisms $\Gamma_i$ of remark (1.1) 
as matrices: 
\[
    \nabla_i(\partial_j) = \sum_k \Gamma_{ij}^k \partial_k.
\]
These symbols completely determine the connection. 

\begin{definition}
    A connection $\nabla$ on $TM$ is called \textit{symmetric} if for any two vector fields $X$ and $Y$
    on $TM$, we have
    \[
        \nabla_X Y - \nabla_Y X = [X, Y].
    \]
\end{definition}

\subsection{Riemannian Geometry}

Suppose $M$ is a Riemannian manifold.

\begin{definition}
    We say $\nabla$ is \textit{compatible with the metric} on $M$ if for any three vector fields
    $X$, $Y_1$, and $Y_2$, we have
    \[
        (\nabla_X Y_1, Y_2) + (Y_1, \nabla_X Y_2) = X \cdot(Y_1, Y_2).
    \]
\end{definition}

\begin{theorem}[Levi-Civita]
    A Riemannian manifold possesses a unique connection that is symmetric and compatible with the metric.
\end{theorem}

\textit{Proof:} Work in local coordinates. Represent the metric by a matrix of smooth functions
$g_{jk} = (\partial_j, \partial_k)$. The symmetry condition implies 
\[
    (\nabla_i \partial_j, \partial_k) + (\partial_j, \nabla_i\partial_k) = \partial_i(\partial_j,\partial_k),
\]
or in other words,
\[
    \partial_i g_{jk} = \sum_a \left(\Gamma_{ij}^a g_{ak} + \Gamma_{ik}^a g_{aj}\right).
\]
Moreover, we can permute indices to get
\[
    \partial_j g_{ki} = \sum_a \left(\Gamma_{jk}^a g_{ai} + \Gamma_{ji}^a g_{ak}\right)\qquad
    \partial_k g_{ij} = \sum_a \left(\Gamma_{ki}^a g_{aj} + \Gamma_{jk}^a g_{ai}\right).
\]

Combining these and using symmetry, we see that 
\[
    \sum_a \Gamma_{ij}^a g_{ak} = \frac{1}{2}\left(\partial_i g_{jk} + \partial_j g_{ik}
    - \partial_k g_{ij}\right).
\]
This determines the $\Gamma$'s uniquely, since $g$ is invertible. \qed 

Since the metric determines the Levi-Civita connection, it also uniquely determines the 
curvature. 

\begin{definition}
    The curvature of the Levi-Civita connection on a Riemannian manifold 
    is called the \textit{Riemann curvature operator} and is denoted by $R$. In components
    relative to a frame $\{e_i\}$ we write 
    \[
        R(e_j, e_k)e_l = \sum_i R_{ljk}^i e_i.
    \]
\end{definition}

We also define the 4-covariant version of the Riemannian curvature tensor defined by 
\[
    R_{iljk} = (R(e_j,e_k)e_l, e_i).
\]

\begin{proposition}
    The Riemann curvature operator has the following properties:
    \begin{enumerate}
        \item[(i)] $R(X,Y)Z + R(Y,X)Z = 0$, or $R_{ljk}^i + R_{lkj}^i = 0$.
        \item[(ii)] $R(X,Y)Z + R(Y,Z)X + R(Z,X)Y = 0$, 
        or $R_{ljk}^i + R_{jkl}^i + R_{klj}^i = 0$.
        \item[(iii)] $(R(X,Y)Z, W) + (R(X,Y)W, Z) = 0$, or $R_{iljk} + R_{lijk} = 0$. 
        \item[(iv)] $(R(X,Y)Z, W) = (R(Z,W)X,Y)$, or $R_{iljk} = R_{jkil}$.
    \end{enumerate}
\end{proposition}

The only non-trivial contraction of the Riemann curvature tensor is called 
the \textit{Ricci curvature tensor}, given by 
\[
    \text{Ric}(Y,Z) \mapsto \text{tr}\left(X \mapsto R(X,Y)Z\right).
\]

In components, $\text{Ric}_{ab} = \sum_i R^i_{aib}$.

We can raise one index of the Ricci curvature tensor to get the Ricci curvature
operator, and taking its trace we obtain the scalar curvature,
\[
    \kappa = g^{ab}\text{Ric}_{ab}.
\]

\begin{definition}
    We say a curve $\gamma$ is a \textit{geodesic} if $\nabla_{\dot{\gamma}}\dot{\gamma} = 0$,
    or in other words if $\dot{\gamma}$ is parallel along $\gamma$. 
\end{definition}

The geodesic equation has a unique solution given initial values $\gamma(0)$ and $\dot{\gamma}(0)$.
Then the exponential map 
\[
    \text{exp}: U \to M
\]
can be defined at any $m \in M$ on some open star-shaped subset $U \subset T_m M$ by 
sending a vector $v \in T_m M$ to the value at $t = 1$ of the unique geodesic 
with starting value $m$ and velocity $v$. By the inverse function theorem, $\text{exp}$
is a diffeomorphism from a neighborhood of zero in $T_m M$ to a neighborhood of $m$ in $M$.
Choosing an orthonormal basis for $T_m M$ gives a coordinate system called
a \textit{geodesic coordinate system} at $m$. 

\begin{proposition}
    At the origin of a geodesic coordinate system, the Christoffel symbols all vanish.
\end{proposition}

\textit{Proof:} It is enough to show that $\nabla_i \partial_j$ vanishes at the origin. 
By symmetry of the connection, it is also enough to show $\nabla_X X = 0$ for
all constant vector fields $X$.

But in a geodesic coordinate system, the geodesics are radial lines from the origin. Moreover,
$X$ clearly points along these geodesics and has constant length, so $\nabla_X X = 0$. \qed

\subsection{Differential forms}

We let $\Omega^p = \bigwedge^p V^*$ denote the space of $p$-forms. We can identify 
$p$-forms with antisymmetric tensors by writing 
\[
    \alpha = \sum_{i_1 < \cdots < i_p} A_{i_1, \dots, i_p} dx^{i_1}\wedge \cdots \wedge dx^{i_p}
    = \frac{1}{p!}\sum_{i_1,\dots, i_p} A_{i_1,\dots, i_p} dx^{i_1}\wedge \cdots \wedge dx^{i_p}.
\]

We first extend the metric to the space of contravariant tensors. 

Let $(\cdot, \cdot)$ be an inner product on $V$, and fix a basis $\{e^i\}$
of $V$, with corresponding dual basis $\{e_i\}$. We have the matrix $g_{ij} = (e_i, e_j)$
and its inverse $g^{ij}$.

Then we have the lowering and raising of incices maps, 
\begin{align*}
    \flat: e^i \mapsto g_{ij}e_j\\
    \sharp: e_i \mapsto g^{ij}e^j.
\end{align*}

These extend to tensors in the natural way.

Using the metric, we can define one on the space of $(0,p)$ tensors by the formula
\[
    (\alpha,\beta) = \frac{1}{k!}\sum_{i_1, \dots, i_p, j_1, \dots, j_p}g^{i_1, j_1}\cdots g^{i_p, j_p}
    A_{i_1, \dots, i_p}B_{i_1, \dots, i_p}.
\]

In an oriented coordinate system, we define the volume form $\text{vol} = \sqrt{g}dx^1 \wedge
\cdots dx^n$, where $g = \det{g_{ij}}$. The normalization constant in the inner product formula ensures that 
$(\text{vol},\text{vol}) = 1$:
\begin{align*}
    (\text{vol}, \text{vol}) = \frac{1}{n!}\sum_{i_1, \dots, i_n, j_1, \dots, j_n}g^{i_1, j_1}\cdots g^{i_n, j_n}
    \delta_{1, \dots, n}^{i_1,\dots, i_n} \delta_{1,\dots,n}^{i_1,\dots, i_n}\\
    = \frac{1}{n!}g \sum_{j_1, \dots, j_p}\det(g^{ij}) = 1.
\end{align*}

We also have formulas for the wedge product and exterior derivative of forms in terms 
of the antisymmetric tensors associated to them. If $C$ corresponds to $\alpha\wedge\beta$, then
\[
    C_{k_1,\dots, k_{p+q}} = \frac{1}{p!q!}\sum_{i_1,\dots, i_p, j_1, \dots, j_q}
    \delta_{k_1, \dots, k_{p+q}}^{i_1,\dots, i_p, j_1,\dots, j_q}.
\]

and if $D \sim d\alpha$, then
\[
    D_{k_1, \dots, k_{p+1}} = \frac{1}{p!}\sum_{i_1, \dots, i_p}\delta^{j,i_1,\dots,i_p}_{k_1,\dots, k_p}
    \frac{\partial A_{i_1,\dots, i_p}}{\partial x_j}.
\]

\begin{definition}
    Let $\alpha$ be a $k$-form. Define the \textit{Hodge Star} $*\alpha$ to be the unique
    $(n-k)$ form such that for all $k$-forms $\beta$, 
    \[
        (\alpha,\beta)\text{vol} = \beta \wedge *\alpha.
    \]
\end{definition}

This operation is linear and we have $**\alpha = (-1)^{kn + k}\alpha$

\begin{definition}
    If $\alpha$ is a $k$-form, define 
    \[
        d^*\alpha = (-1)^{nk + n + 1}*d*\alpha.
    \]
    Then $d^*\alpha$ is a $(k-1)$ form, and $(d^*)^2 = 0$.
\end{definition}

The importance of $d^*$ is that it is the adjoint of $d$, once we define a global inner product 
on forms.

We define
\[
    \langle \alpha, \beta\rangle = \int_M (\alpha,\beta)\text{vol} = 
    \int_M \beta \wedge *\alpha = \int_M 
\]

\subsection{Excercises}

\begin{enumerate}
    \item[1.28] Prove Cartan's formula for the exterior deriative of a $p$-form.
    
    \textit{Proof:} See homework either 6 or 7 from DG.

    \item[1.29] Prove the second Bianchi Identity,
    \[
        \nabla R(X,Y,Z,V,W) + \nabla R(X,Y,V,W,Z) + \nabla R(X,Y,W,Z,V) = 0.
    \]

    \textit{Proof:} Since the objects are tensors, it suffices to prove 
    the identity at an arbitrary point on $M$. Fix a geodesic coordinate
    system about this point so that brackets and $\nabla_i e_j$ vanishes 
    at the origin. 

    Then follow the steps in Do Carmo with the simplified coordinate system,
    then the result is clear.

    \item[1.30] Let $x^i$ be a geodesic coordinate system on a manifold $M$, and
    let $r$ denote the Riemannian distance from the point $x$ to the origin. Prove
    that 
    \[
        \sum_{i,j} g_{ij}(x)x^i x^j = \sum_{i,j} g^{ij}(x)x^i x^j = r^2.
    \]

    \textit{Proof:} This problem is straightforward using the 
    fact that in a geodesic coordinate system, the geodesics are just radial
    lines from the origin. The distance from the origin to $x$ is the infimum over
    all geodesics connecting these points of the length of these curves. But 
    in this case, there is only one radial line, so its length is the inf. 

    Moreover, the velocity of the geodesic at the origin at $t = 0$ and at $x$ at $t = 1$ 
    is just the magnitude of the vector $x$, viewed as a subset of $T_0 M$. 

    \item[1.31] Let $\alpha$ be a $k$-form on an $n$ -dimensional oriented Riemannian 
    manifold $M$ represented in local coordinates by an antisymetric tensor $A$. 
    Show that 
    \[
        *\alpha = \sum_{i_1, \dots, i_k, j_1, \dots, j_{n-k}} 
        \frac{1}{(n-k)!k!}\delta^{1, \dots, n}_{i_1, \dots, i_k, j_1, \dots, j_{n-k}}
        \sqrt{g} A^{i_1, \dots, i_k}dx^{j_1} \wedge \cdots \wedge dx^{j_{n-k}}.
    \]

    \textit{Proof:} We utilize the antisymmetric tensor representation formula for the wedge
    product:
    \[
        C_{k_1,\dots, k_{p+q}} = \frac{1}{p!q!}\sum_{i_1,\dots, i_p, j_1, \dots, j_q}
        \delta_{k_1, \dots, k_{p+q}}^{i_1,\dots, i_p, j_1,\dots, j_q},
    \]
    where $C$ is the tensor representing the wedge of a $p$-form and a $q$-form.

    We also have a formula for the inner product of $k$ forms:
    \[
        (\alpha,\beta) = \frac{1}{k!}\sum_{i_1, \dots, i_p, j_1, \dots, j_p}g^{i_1, j_1}\cdots g^{i_p, j_p}
        A_{i_1, \dots, i_p}B_{i_1, \dots, i_p}.
    \]

    By definition of $*\alpha$, for all $k$-forms $\beta$, we have
    \[
        (\alpha,\beta)\text{vol} = \beta \wedge *\alpha.
    \]
    In other words, if we write $*A$ and $B$ for the antisymmetric tensors corresponding to 
    $*\alpha$ and $\beta$, this means
    \[
        \sqrt{g}\frac{1}{k!}\sum_{|I|=|J|=k}g^{I,J}A_I B_J 
        = \frac{1}{(n-k)!k!}\sum_{|I|=k,|J|=n-k}\delta^{IJ}_{1,\dots, n}B_I (*A)_J.
    \]

    We start from the left. Notice this is just the raising of the indices of $A$:
    \[
        A^J = \sum_{|I| = |J| = k}g^{I,J}A_I.
    \]
    Then we have 
    \[
        (\alpha,\beta)\text{vol} = \sqrt{g}\frac{1}{k!}\sum_{|J| = k}A^J B_J.
    \]

    Now we need a trick with the generalized Kronoker-delta symbols. For 
    $k$-tuples $I,J$ and an $n-k$-tuple $K$, we have
    \[
        \delta^{IK}_{1,\dots,n}\delta^{JK}_{1,\dots,n}\delta^I_J = 1.
    \]
    To see why, consider either when $I$ and $J$ are opposite sign, 
    or the same.

    Then we can write
    \begin{align*}
        \sqrt{g}\frac{1}{k}\sum_{|J| = k}A^J B_J
        & = \frac{1}{(n-k)!k!}\sqrt{g}\frac{1}{k!}
        \sum_{|I|=|J|=k, |K|=n-k}\delta^{IK}_{1,\dots,n}\delta^{JK}_{1,\dots,n}\delta^I_J
        B_J A^J\\
        & = \frac{1}{(n-k)!k!}\sqrt{g}\frac{1}{k!} \sum_{|I|=|J|=k, |K|=n-k}
        \delta^{IK}_{1,\dots,n}\delta^{JK}_{1,\dots,n}B_J A^I\\
        & = \frac{1}{(n-k)!k!} \sum_{|J|=k, |K|= n-k}\delta^{IJ}_{1,\dots,n} B_I
        \left(\sqrt{g}\frac{1}{k!}\sum_{|I|=k}\delta^{IK}_{1,\dots,n}A^I\right)\\
        & = B\wedge *A,
    \end{align*}

    where we write $B \wedge *A$ to be the tensor for $\beta\wedge *\alpha$. In other words,
    the matrix in parenthesis with parameter $K$ is $*A$. 
\end{enumerate}

\section{The Spin Groups}

\begin{definition}
    An algebra $A$ over $\RR$ or $\CC$ is called a \textit{superalgebra} 
    if it is an internal direct sum of linear subspaces $A_0$ and $A_1$, 
    with
    \[
        A_0\cdot A_0 \subset A_0, \quad A_0 \cdot A_1 \subset A_1, \quad 
        A_1\cdot A_0 \subset A_1, \quad A_1 \cdot A_1 \subset A_0.
    \]
    We say that $A_0$ and $A_1$ are the \textit{even} and \textit{odd} parts of $A$
    respectively. The subset $A_0\cup A_1$ is called the set of homogeneous elements of $A$. 

    Alternatively, a superalgebra is an algebra $A$ equipped with an automorphism 
    $\epsilon$, the \textit{grading automorphism}, such that $\epsilon^2 = 1$. 
    The subspaces $A_0$ and $A_1$ are the $1$ and $-1$ eigenspaces of $\epsilon$, 
    and in the previous case the automorphism is defined by
    \[
        \epsilon(a_0 + a_1) = a_0 - a_1
    \]
    for $a_0 \in A_0, a_1 \in A_1$. 
\end{definition}

\begin{proposition}
    The Clifford algebra of an inner product space $V$ is a super algebra, in which the homogeneous
    elements are products of an even or odd number of elements of $V$. 
\end{proposition}

\begin{definition}
    We define the \textit{supercommutator} on $A$ first on homogeneous elements by
    \[
        [x,y]_s = xy - (-1)^{\deg(x)\deg(y)}yx,
    \]
    and then extend to all of $A$ by linearity. We define the \textit{super-center}
    $\mathfrak{Z}_s(A)$ of $A$ to be the set of all elements which \textit{supercommute}
    with every element in $A$, i.e. $[x,y]_s = 0$.
\end{definition}

\end{document}